\documentclass[11pt,a4paper]{article}
\usepackage[margin=2.5cm]{geometry}
\usepackage[T1]{fontenc}
\usepackage[utf8]{inputenc}
\usepackage{booktabs}
\usepackage{enumitem}
\usepackage{hyperref}
\usepackage{parskip}

% Rupee symbol: prints "Rs." for maximum compiler compatibility.
% If compiling with XeLaTeX/LuaLaTeX and a suitable font, replace with ₹.
\newcommand{\rupee}{Rs.\,}

\title{Varsap: A Zero-Cost WhatsApp Feedback Automation Platform\\
\large Design, Implementation and Production Deployment of a Customer
Feedback Engine for an Amazon Seller}
\author{Sathvik Kilari}
\date{July 2026}

\begin{document}
\maketitle

\begin{abstract}
This report presents Varsap, a full-stack web application that automates
WhatsApp-based customer feedback collection for Varistor, an Amazon
marketplace seller. The system ingests Amazon order exports, dispatches
personalised WhatsApp template messages through Meta's WhatsApp Cloud API,
and tracks each message through its complete delivery lifecycle
(sent~$\rightarrow$~delivered~$\rightarrow$~read) in real time. A defining
engineering constraint was a total infrastructure budget of
\textbf{\rupee 0 per month}, satisfied by composing free-tier cloud
services: serverless hosting, a managed document database, and a managed
message queue. The project covers the complete software lifecycle ---
requirements, architecture, implementation, production debugging against a
third-party platform, security hardening, cost engineering, and a formal
handover to the company, including migration of all platform assets to
corporate ownership. The system is live in production and has been verified
end-to-end with real customer traffic.
\end{abstract}

\section{Introduction and Problem Statement}

Amazon sellers depend on product reviews and buyer feedback, yet the
marketplace itself offers sellers no direct, high-engagement channel to
request it. Email requests suffer from low open rates, while WhatsApp ---
the dominant messaging platform among Indian consumers --- offers open
rates above 90\%, but is aggressively protected against spam: businesses
may only initiate conversations using message templates pre-approved by
Meta, sent via the paid WhatsApp Business (Cloud) API.

The goal of this project was to give Varistor's non-technical staff a
one-click pipeline from ``here is today's Amazon order export'' to
``every buyer has received a personalised feedback request on WhatsApp,
and we can watch delivery and read receipts live'' --- subject to three
hard requirements:

\begin{enumerate}[nosep]
  \item \textbf{Zero fixed cost}: no paid infrastructure of any kind; the
        only permissible spend is Meta's unavoidable per-message fee.
  \item \textbf{Operable by non-technical employees}: no console access,
        no configuration files, no developer in the loop for daily use.
  \item \textbf{Safe by construction}: no customer may ever be messaged
        twice for the same order, opt-outs must be honoured permanently,
        and access must be restricted to authorised staff.
\end{enumerate}

\section{System Architecture}

\subsection{Technology stack}

\begin{table}[h]
\centering
\begin{tabular}{lll}
\toprule
\textbf{Layer} & \textbf{Technology} & \textbf{Free-tier role} \\
\midrule
Web application & Next.js (App Router, TypeScript) & --- \\
Hosting \& CI/CD & Vercel & Auto-deploy on git push \\
Database & MongoDB Atlas (M0) & Campaigns, messages, blocklist \\
Message queue & Upstash QStash & Paced delivery, retries, scheduling \\
Messaging & Meta WhatsApp Cloud API & Template send + status webhooks \\
Authentication & NextAuth + Google OAuth & Domain-restricted SSO \\
\bottomrule
\end{tabular}
\caption{Composition of free-tier services meeting the \rupee 0/month constraint.}
\end{table}

\subsection{Campaign data flow}

A campaign proceeds through four fully decoupled stages:

\begin{enumerate}[nosep]
  \item \textbf{Ingest} --- the employee uploads an Excel/CSV export;
        columns are mapped interactively with a live preview. Rows are
        validated (phone normalisation, order-ID presence) client-side.
  \item \textbf{Create} --- the server deduplicates against all prior
        campaigns using a unique database index on the
        \emph{(phone, order-ID)} pair, removes blocklisted customers, and
        presents a pre-flight cost estimate before anything is sent.
  \item \textbf{Dispatch} --- every message is handed to the QStash queue
        with a per-message timestamp (one message per second), so sending
        proceeds serverlessly in the background; the employee may close
        the browser immediately. Campaigns may also be scheduled for a
        future time.
  \item \textbf{Track} --- Meta pushes signed webhook events (sent,
        delivered, read, failed, and inbound replies) to the application,
        which updates the dashboard in real time. A customer replying
        ``STOP'' is inserted into a permanent blocklist and all their
        in-flight messages are cancelled.
\end{enumerate}

A key architectural property is that \textbf{no sending logic runs in the
browser}: the queue drives delivery and the webhook maintains state,
making the system resilient to the employee's machine, connection, or
session.

\subsection{Reliability and safety mechanisms}

\begin{itemize}[nosep]
  \item \emph{Idempotent sending}: queue retries can never double-send; a
        message is only dispatched from the \texttt{queued} state.
  \item \emph{Monotonic status lifecycle}: a late ``delivered'' event can
        never downgrade an already-``read'' message.
  \item \emph{Dedupe by design}: the unique \emph{(phone, order-ID)} index
        makes re-uploading the same file harmless --- duplicates are
        skipped and reported, not resent.
  \item \emph{Fail-closed webhook security}: every webhook call is
        HMAC-SHA256 signature-verified against the Meta app secret using
        constant-time comparison; an unconfigured secret rejects requests
        rather than silently accepting unsigned ones.
  \item \emph{Machine-endpoint authentication}: the queue worker verifies
        QStash request signatures; all human-facing routes require a
        whitelisted Google login (company domain plus an explicit
        allowlist).
\end{itemize}

\section{Feature Set}

\begin{itemize}[nosep]
  \item Campaign creation from arbitrary Amazon exports (CSV/XLSX) with
        interactive column mapping and recipient preview.
  \item Pre-flight cost estimation and post-dedupe final cost display.
  \item Immediate or scheduled dispatch; safe retry of failed messages only.
  \item Live per-message dashboard: pending/queued/sent/delivered/read/%
        failed/opted-out, with failure reasons surfaced from Meta.
  \item Automatic opt-out (STOP) handling with permanent blocklisting.
  \item \textbf{In-app template management}: employees create WhatsApp
        message templates and submit them to Meta for approval from within
        the application --- no Meta console access needed --- and track
        approval status live.
  \item \textbf{Per-campaign template picker}: each campaign selects any
        approved template from a dropdown, with a WhatsApp-style preview
        of the exact message (variables filled with sample values) so
        non-technical staff see precisely what customers will receive.
  \item Duplicate-upload warning heuristic (same row count within 24h)
        layered on top of the hard database guarantee.
\end{itemize}

\section{Production Debugging Case Studies}

Integrating against a large third-party platform produced several
non-trivial failures whose diagnosis forms a significant part of the
project's engineering contribution.

\subsection{Case 1: Silent loss of delivery statuses}
\textbf{Symptom:} messages delivered on phones, but the dashboard froze at
``sent''. \textbf{Root cause:} the webhook signature secret had never been
configured in production; the deliberately fail-closed design was
rejecting every status callback. \textbf{Resolution:} configure the app
secret and verify with Meta's webhook self-test plus a live inbound
message. \textbf{Lesson:} fail-closed security is correct, but must be
paired with observability --- the failure was invisible until inspected at
the platform logs level.

\subsection{Case 2: Systematic template rejection (\texttt{INVALID\_FORMAT})}
\textbf{Symptom:} every template submitted through the application was
instantly rejected by Meta, while seemingly identical templates created
earlier were approved. \textbf{Investigation:} the Graph API's
\texttt{rejected\_reason} field, per-template \texttt{components}
introspection, and controlled resubmission experiments isolated the
variable. \textbf{Root cause:} Meta's template API defaults to
\emph{positional} parameters (\texttt{\{\{1\}\}}); the application uses
\emph{named} parameters (\texttt{\{\{name\}\}}) but did not declare
\texttt{parameter\_format: NAMED}, so Meta parsed the named variable as
malformed, silently discarded the required example values, and
auto-rejected. A one-line fix in the template-creation payload resolved
it permanently. Secondary findings: Meta auto-rejects near-duplicate
wording of existing templates, and rejected/deleted template names are
locked for roughly four weeks. \textbf{Lesson:} when a black-box reviewer
rejects input, bisect by reproducing via the raw API and diff what the
platform \emph{stored} against what was \emph{sent}.

\subsection{Case 3: Credential lifetime engineering}
\textbf{Symptom:} recurring authentication failures --- developer-console
access tokens expire every 24 hours, and one expiry silently broke local
development while production continued on a separately-updated token.
\textbf{Resolution:} replaced ad-hoc tokens with a Meta \emph{system user}
(machine identity) holding full access to the app and both WhatsApp
accounts, issuing a non-expiring token verified via the
\texttt{debug\_token} introspection endpoint. \textbf{Lesson:} humans
should authenticate interactively; systems should authenticate with
machine identities whose lifetime is an explicit, audited decision.

\section{Security Model}

\begin{itemize}[nosep]
  \item All secrets exist only in the hosting platform's environment
        configuration; the repository contains placeholders only.
  \item Authentication: Google SSO restricted to the company domain plus
        an explicit allowlist; every page and API route is gated except
        the public privacy page and the two signature-verified machine
        endpoints.
  \item Webhook integrity: HMAC-SHA256 with constant-time comparison,
        failing closed.
  \item Credential hygiene: non-expiring machine token with documented
        invalidation conditions; a written operational runbook forbids the
        four actions that would sever credentials (deleting the system
        user, resetting the app secret, revoking tokens, transferring
        assets between business portfolios).
\end{itemize}

\section{Cost Engineering}

Fixed infrastructure cost is \rupee 0/month by construction. Variable cost
is Meta's per-message fee, which in India differs roughly sevenfold by
template category: utility ($\approx$\rupee 0.12/message) versus marketing
($\approx$\rupee 0.88/message). The application makes this visible at the
point of decision: the template picker labels each template's category,
the pre-flight estimator prices the campaign before creation, and the
feedback template is deliberately maintained in the cheaper utility
category where policy allows. Customer replies are free and open a
24-hour service window, which the message copy deliberately encourages
(``just reply to this message'').

\section{Deployment, Handover and Asset Migration}

The system was brought to production on the company's WhatsApp number
(display name ``Varistor'', quality rating green) with Meta webhooks,
billing, and a published Meta app. The final project phase was
organisational rather than technical: transferring ownership so the
company controls all keys and billing. Deliverables include:

\begin{itemize}[nosep]
  \item \texttt{KT.md} --- a complete knowledge-transfer document:
        architecture, environment variables, operating guide for daily
        staff, troubleshooting matrix, and a handover checklist.
  \item \texttt{MIGRATION.md} --- a phased runbook for migrating the Meta
        app (official app transfer), the WhatsApp number (official number
        migration between business accounts), templates, billing, and
        machine credentials to the company's own verified Meta Business
        Manager, with explicit verification gates before each destructive
        step.
  \item An end-to-end verification protocol (fresh-order test campaign
        walking sent~$\rightarrow$~delivered~$\rightarrow$~read) used as
        the acceptance criterion for every configuration change.
\end{itemize}

\section{Results}

\begin{itemize}[nosep]
  \item Live production system, verified end-to-end with real WhatsApp
        traffic on the company's number, at \rupee 0/month fixed cost.
  \item Full message-lifecycle tracking with sub-minute latency from
        WhatsApp read receipt to dashboard update.
  \item Template management and per-campaign template selection usable by
        non-technical staff, including live message previews.
  \item Three production incidents root-caused and permanently fixed, each
        contributing a documented, transferable lesson.
  \item Complete, documented corporate handover path with credentials
        engineered to require zero ongoing maintenance.
\end{itemize}

\section{Limitations and Future Work}

\begin{itemize}[nosep]
  \item \textbf{No reply inbox}: inbound customer replies are stored for
        opt-out detection but have no reading UI; feedback currently
        arrives on the business phone. An inbox view is the natural next
        feature.
  \item \textbf{Category-aware cost estimation}: the estimator uses a
        single configured rate; wiring it to the selected template's
        category would price marketing campaigns accurately.
  \item \textbf{Throughput ceilings}: the free queue tier ($\approx$500
        messages/day) and Meta's pre-verification conversation limit
        (250/day, lifted by business verification) bound current scale;
        both have documented upgrade paths.
  \item \textbf{Analytics}: read-rate and reply-rate aggregation across
        campaigns would close the feedback loop on message copy quality.
\end{itemize}

\section{Learning Outcomes}

The project exercised the full lifecycle of a production integration:
translating a business need into hard constraints; composing free-tier
managed services into a reliable architecture; defensive design
(idempotency, dedupe-by-index, fail-closed verification); black-box
debugging of a third-party review system via API introspection and
controlled experiments; credential and identity engineering (machine
identities, token introspection, documented invalidation conditions);
cost-aware product design; and the organisational discipline of
documentation, knowledge transfer, and asset migration with verification
gates. The overarching lesson: in production systems, the code is the
smaller half --- the platform's rules, the credentials' lifetimes, and
the operators' understanding are what keep a system alive.

\end{document}
